% SPDX-FileCopyrightText: 2026 Yiyu Zhou <yiyu@yiyuzhou.io>
%
% SPDX-License-Identifier: GPL-3.0-or-later

\chapter{Fading in and out}
\label{ch:fades}

A fade cue changes the level of another cue, its target, over time.  This
chapter runs the fades of the storm show, then adds one to your own show.

\section{Fades in the storm}

Press \key{Control+l} and open \texttt{storm.cinq.toml}.

\begin{enumerate}
\item Select cue~5, Rain, press \ui{Standby selected}, and press \ui{GO}\@.
  The rain starts quietly, at $-20$~dB, and loops.
\item Press \ui{GO} again to fire cue~6, Rain up.  Over eight seconds, the
  rain rises to $-6$~dB and keeps playing.
\item Press \ui{GO} twice more.  Cue~7 plays a thunder crack, and cue~8,
  Rain out, fades the rain out over ten seconds and stops it.
\end{enumerate}

Select cue~6 to see how it works.  On the \ui{Basics} tab, its \ui{Target}
is cue~5 and its \ui{Duration} is 8 seconds.  On the \ui{Levels} tab, shown
in Figure~\ref{fig:levels}, \ui{Fade target to} sets the level at which the
fade ends.  \ui{Stop target when done} is cleared, so the rain keeps playing
until cue~8 fades it out.

\begin{figure}[H]
  \centering
  \includegraphics[width=\textwidth]{screenshots/levels.png}
  \caption{\ui{Levels} tab of a fade cue that leaves its target playing}
  \label{fig:levels}
\end{figure}

You can change levels while cues play, which helps in a rehearsal.  Change
the \ui{Fade target to} of a fade cue while it fades, and the fade turns
toward the new level in the time it has left.  Change the level of the rain
itself after a fade has lowered it, and the rain moves by the same amount,
so that a cue that a fade has lowered stays lower.

A check box, such as \ui{Stop target when done}, takes the keyboard with
\key{Tab} like a field, and \key{Enter} checks or clears it.

\section{A fade of your own}

Press \key{Control+l} and open \texttt{my-show.cinq.toml}.

\begin{enumerate}
\item Select the rain, choose \ui{Fade} in the \ui{New cue} drop-down list of
  the toolbar, and press \ui{Add cue}.  The new fade cue appears below the
  rain and targets it, since a new fade cue targets the cue above it.  By
  default, it fades its target to silence over five seconds and then stops
  it.  With the keys, press \key{Shift+Tab} twice from the cue list, which
  reaches the \ui{New cue} list from the end of the order of the fields,
  since \key{Tab} would pass the fields of the inspector first, press
  \key{Enter} to open the list, move to \ui{Fade} with \key{j} and \key{k},
  and press \key{Enter} to pick it, which gives the keyboard back to the cue
  list.  Then press \key{o}, which adds a cue of the kind that the list shows.
  The list keeps its kind, so each later \key{o} adds a fade cue until you
  pick another kind.
\item Select the rain, press \ui{Standby selected}, and press \ui{GO} twice:
  once to start the rain, and once to fade it out.
\end{enumerate}

The \ui{Target} list on the \ui{Basics} tab selects a different target, and
the \ui{Duration} field sets a different length.  Clear \ui{Stop target when
done} only for a fade that should leave its target playing, as cue~6 of the
storm show does; a cue faded to silence but not stopped keeps running,
inaudibly.  A stop cue, by contrast, stops its target immediately, without a
fade.

\section*{Exercise}

You now know enough to build a show of sounds: cues that play them, and fades
that bring them up and take them out.  Add a fade cue that brings your rain
in from silence: set the rain's level to $-60$~dB, and fade it to $-10$~dB
without stopping it.
